\chapter{Frozen Shoulder Release}

\section*{What Is This Surgery?}
Frozen shoulder release, also known as shoulder capsular release, is a surgical intervention aimed at alleviating the debilitating pain and restricted range of motion associated with adhesive capsulitis, commonly referred to as frozen shoulder. This condition is characterized by a progressive thickening, inflammation, and contraction of the glenohumeral joint capsule, leading to the formation of adhesions that severely limit both active and passive shoulder movement. The surgical procedure involves mechanically dividing these stiffened and contracted capsular tissues, particularly in the rotator interval and the anterior and inferior aspects of the capsule, thereby restoring the normal volume and mobility of the shoulder joint. Typically performed arthroscopically, this surgery often includes a controlled manipulation of the arm under anesthesia to achieve immediate and significant gains in shoulder mobility. The ultimate goal is to enable patients to resume their daily activities, work, and recreational pursuits without the debilitating pain and stiffness that characterize this condition.

\section*{Alternative Names}
\begin{itemize}
  \item Arthroscopic Capsular Release
  \item Shoulder Arthrolysis
  \item Manipulation Under Anesthesia (MUA) for Frozen Shoulder
  \item Arthroscopic Capsulotomy
  \item Open Capsular Release (less common)
  \item Surgical Release for Adhesive Capsulitis
\end{itemize}

\section*{Relevant Surgical Anatomy}
The glenohumeral joint is a ball-and-socket synovial joint formed by the articulation of the humeral head with the glenoid fossa of the scapula. The joint capsule, a fibrous envelope, is crucial for stability and range of motion. In adhesive capsulitis, the anterior, inferior, and posterior aspects of this capsule, along with the rotator interval (the space between the supraspinatus and subscapularis tendons), become pathologically thickened, fibrotic, and contracted. Key structures involved in the pathology and targeted during surgery include the coracohumeral ligament, superior glenohumeral ligament, and the anterior band of the inferior glenohumeral ligament complex.

The surrounding neurovascular structures are critical considerations. The axillary nerve, which innervates the deltoid and teres minor muscles, courses inferiorly to the glenohumeral joint and is particularly vulnerable during inferior capsular release or forceful manipulation. The musculocutaneous nerve, supplying the biceps and brachialis, and the suprascapular nerve, supplying the supraspinatus and infraspinatus, are also in proximity. The arterial supply to the shoulder region primarily comes from the anterior and posterior circumflex humeral arteries, which are branches of the axillary artery, and the suprascapular artery. Venous drainage largely parallels the arterial supply. Lymphatic drainage from the shoulder region typically flows into the axillary lymph nodes. Surgical landmarks include the acromion, coracoid process, and the palpable borders of the deltoid and pectoralis major muscles, which guide portal placement for arthroscopy.

\section*{Surgical Principle \& Rationale}
The fundamental principle of frozen shoulder release is the mechanical division of the pathologically contracted and fibrotic glenohumeral joint capsule and associated ligaments to restore the normal volume and mobility of the joint. In adhesive capsulitis, the capsule undergoes inflammatory and fibrotic changes, leading to a significant reduction in joint space and adherence of the capsular folds, particularly in the rotator interval, anterior, and inferior aspects. This restricts the normal gliding and rotation of the humeral head within the glenoid fossa.

The rationale for surgical intervention arises when conservative treatments, such as physiotherapy, oral medications, and corticosteroid injections, fail to provide adequate relief or restore functional range of motion over an extended period. The operative concept involves using arthroscopic instruments to precisely incise and release the thickened capsule, the coracohumeral ligament, and the rotator interval under direct visualization. This process effectively "unfreezes" the joint, allowing for immediate gains in passive range of motion. Often, a gentle manipulation under anesthesia is performed after the capsular release to ensure all adhesions are broken and to confirm full glenohumeral excursion. The technical goals are to achieve a comprehensive release without causing iatrogenic injury to surrounding neurovascular structures or articular cartilage, thereby facilitating an accelerated and more complete rehabilitation process.

\section*{History \& Surgical Evolution}
The treatment of frozen shoulder has evolved significantly over centuries. Early interventions for joint stiffness often involved forceful manipulation, sometimes without anesthesia, leading to unpredictable outcomes and frequent complications. The concept of manipulation under anesthesia (MUA) gained prominence in the early 20th century, offering a more controlled method to break adhesions and restore motion. However, MUA alone carried risks of humeral fracture, rotator cuff tears, and nerve injury, especially with excessive force.

The mid-to-late 20th century witnessed the advent of arthroscopy, revolutionizing shoulder surgery. Dr. Lanny Johnson is credited with pioneering arthroscopic techniques in the 1970s. By the 1980s and 1990s, arthroscopic capsular release emerged as a safer and more precise alternative to blind MUA or open capsulotomy. Surgeons like Dr. Stephen Snyder and Dr. Eugene Wolf contributed to refining arthroscopic techniques, particularly in identifying and releasing the critical structures such as the rotator interval and the anterior and inferior capsule. Modern practice often combines a meticulous arthroscopic release with a gentle, controlled MUA to maximize motion gains while minimizing risks. The emphasis has also shifted towards early and aggressive postoperative physiotherapy, recognized as crucial for maintaining the surgically achieved motion and preventing recurrence of stiffness.

\section*{Clinical Indications}
Frozen shoulder release is typically considered when patients meet specific criteria, indicating a failure of conservative management and significant functional impairment.

\begin{itemize}
  \item Persistent, severe pain and restricted range of motion in adhesive capsulitis that has failed to improve after a minimum of 3 to 6 months of structured, intensive nonoperative treatment, including physiotherapy, home exercise programs, and nonsteroidal anti-inflammatory medications.
  \item Significant limitation of activities of daily living (ADLs), work, or recreational activities due to shoulder stiffness and pain, despite adequate pain control with oral analgesics or intra-articular corticosteroid injections.
  \item Documented failure of intra-articular corticosteroid injections to provide sustained relief or improve range of motion.
  \item Patients in the "frozen" or "thawing" phase of adhesive capsulitis who desire a more rapid restoration of function and are unable to tolerate the prolonged natural history of the condition.
  \item Secondary adhesive capsulitis associated with systemic conditions such as diabetes mellitus or thyroid disease, where the stiffness is particularly recalcitrant to conservative measures and medical management of the underlying condition is optimized.
  \item Objective findings of severely restricted passive range of motion in multiple planes, particularly external rotation and abduction, confirmed by physical examination.
  \item Absence of other primary shoulder pathologies, such as significant rotator cuff tears, advanced glenohumeral osteoarthritis, or calcific tendinitis, that would contraindicate or limit the efficacy of capsular release.
  \item A desire to accelerate recovery when prolonged conservative management is impractical due to occupational demands or patient preference, after thorough discussion of risks and benefits.
\end{itemize}

\section*{Clinical Positioning}
Frozen shoulder release is positioned as a secondary or salvage treatment option within a stepped care approach for adhesive capsulitis. It is generally not considered a primary intervention due to the self-limiting nature of the condition in many patients and the effectiveness of nonoperative management. The vast majority of patients with frozen shoulder will initially undergo a comprehensive course of conservative therapy, which typically includes a combination of physical therapy focusing on stretching and range-of-motion exercises, oral analgesics, nonsteroidal anti-inflammatory drugs (NSAIDs), and often one or more intra-articular corticosteroid injections.

Surgery is reserved for those patients who experience persistent, debilitating pain and stiffness that significantly impairs their quality of life and functional capacity, despite at least three to six months of diligent and well-structured conservative treatment. In these refractory cases, surgical release becomes the definitive step to mechanically break the cycle of inflammation and fibrosis, offering a more rapid and predictable restoration of motion compared to continued conservative management. For patients with secondary adhesive capsulitis, particularly those with diabetes, the condition can be more severe and prolonged, making surgical intervention a more frequent consideration when conservative measures fail. Therefore, while not a first-line treatment, surgical release plays a crucial role as an effective intervention for carefully selected patients who have exhausted nonoperative pathways.

\section*{Surgical Technique \& Procedure}
Frozen shoulder release is typically performed under general anesthesia, often supplemented with an interscalene regional nerve block to provide excellent postoperative pain control and facilitate early rehabilitation. The patient is usually positioned in the lateral decubitus position with the affected arm draped free, or in a beach-chair position, allowing for full manipulation of the shoulder.

The procedure generally involves the following steps:

\begin{itemize}
  \item \textbf{Diagnostic Arthroscopy:} Standard posterior and anterior portals are established. A 30-degree arthroscope is introduced through the posterior portal to perform a thorough diagnostic assessment of the glenohumeral joint, confirming the diagnosis of adhesive capsulitis and ruling out other intra-articular pathologies. The characteristic findings include a globally contracted joint capsule, obliteration of the rotator interval, and often hypertrophic synovitis.
  
  \item \textbf{Rotator Interval Release:} Using an arthroscopic shaver or radiofrequency ablator, the contracted rotator interval, including the superior glenohumeral ligament and the coracohumeral ligament, is meticulously released. This is a critical step for restoring external rotation.
  
  \item \textbf{Anterior Capsular Release:} The anterior capsule, particularly the middle and inferior glenohumeral ligaments, is released from the glenoid rim. Care is taken to protect the axillary nerve, which lies inferiorly.
  
  \item \textbf{Inferior Capsular Release:} The inferior capsule is released from its attachment to the glenoid and humeral neck. This is often the most contracted portion and is crucial for restoring abduction and external rotation. Extreme caution is exercised to avoid injury to the axillary nerve, which is in close proximity.
  
  \item \textbf{Posterior Capsular Release (as needed):} If significant posterior stiffness persists, the posterior capsule may also be released, though this is less commonly the primary target.
  
  \item \textbf{Manipulation Under Anesthesia (MUA):} After the capsular release, the surgeon gently and systematically manipulates the arm through its full range of motion (flexion, extension, abduction, adduction, internal and external rotation) to ensure all remaining adhesions are broken and to confirm a free glenohumeral excursion. This step is performed with controlled force to prevent iatrogenic fracture or nerve injury.
  
  \item \textbf{Joint Lavage and Closure:} The joint is thoroughly irrigated, and the arthroscopic portals are closed with one or two sutures or adhesive strips. A sterile dressing is applied.
\end{itemize}

The procedure typically lasts 30 to 60 minutes and is usually performed as an outpatient day surgery. Drains are rarely used. Immediate postoperative physiotherapy is initiated to maintain the motion gained during surgery.

\section*{Preoperative Workup \& Preparation}
A thorough preoperative workup is essential to ensure patient safety and optimize surgical outcomes for frozen shoulder release. This typically begins with a detailed medical history and physical examination, focusing on the duration and severity of symptoms, previous treatments, and any comorbidities.

\begin{itemize}
  \item \textbf{Medical Evaluation:} Patients undergo a comprehensive medical evaluation by their primary care physician or an internist to assess overall health, particularly cardiac and pulmonary status, and to optimize management of chronic conditions such as diabetes, hypertension, or thyroid disease. Preoperative laboratory tests, including a complete blood count, electrolyte panel, renal and liver function tests, and coagulation profile, are obtained. For diabetic patients, strict glycemic control (HbA1c < 7.0\%) is crucial to minimize infection risk and improve healing.
  
  \item \textbf{Imaging Studies:} Magnetic resonance imaging (MRI) of the shoulder is often performed to confirm the diagnosis of adhesive capsulitis, rule out other causes of shoulder pain and stiffness (e.g., rotator cuff tears, labral tears, advanced osteoarthritis, calcific tendinitis), and assess the extent of capsular thickening. X-rays are also obtained to evaluate for arthritis or previous fractures.
  
  \item \textbf{Medication Review:} All medications are reviewed. Anticoagulants (e.g., warfarin, direct oral anticoagulants) and antiplatelet agents (e.g., aspirin, clopidogrel) are typically held for a specified period before surgery, as per institutional guidelines and in consultation with the prescribing physician, to minimize bleeding risk.
  
  \item \textbf{NPO Status:} Patients are instructed to be NPO (nil per os) after midnight on the day of surgery to prevent aspiration during anesthesia.
  
  \item \textbf{Prophylactic Antibiotics:} Intravenous prophylactic antibiotics, typically a first- or second-generation cephalosporin (e.g., cefazolin), are administered within 60 minutes prior to incision to reduce the risk of surgical site infection.
  
  \item \textbf{Informed Consent:} A detailed discussion with the patient regarding the surgical procedure, potential benefits, risks, alternatives, and the critical role of postoperative rehabilitation is conducted, and informed consent is obtained.
  
  \item \textbf{Postoperative Planning:} Patients are advised to arrange for transportation home and to schedule their initial physiotherapy appointments immediately following surgery, as early and consistent rehabilitation is paramount for a successful outcome.
\end{itemize}

\section*{Contraindications \& Precautions}
Careful patient selection is crucial for successful frozen shoulder release. Both absolute and relative contraindications must be considered.

\textbf{Absolute Contraindications:}
\begin{itemize}
  \item Active infection within the shoulder joint (septic arthritis) or overlying skin (cellulitis), as surgery could spread the infection.
  
  \item Uncontrolled bleeding disorder or coagulopathy that cannot be safely corrected preoperatively.
  
  \item Severe, uncompensated medical comorbidities (e.g., unstable angina, recent myocardial infarction, severe uncontrolled heart failure, severe respiratory insufficiency) that pose an unacceptably high anesthetic or surgical risk.
  
  \item Patient refusal or inability to comply with the intensive postoperative rehabilitation protocol.
\end{itemize}

\textbf{Relative Contraindications \& Precautions:}
\begin{itemize}
  \item \textbf{Severe Osteoporosis:} Increases the risk of iatrogenic humeral fracture during manipulation under anesthesia. In such cases, manipulation should be performed with extreme gentleness, or an arthroscopic release without manipulation may be preferred.
  
  \item \textbf{Recent Proximal Humerus or Glenoid Fracture:} Manipulation could disrupt healing bone. Surgery should be delayed until adequate fracture healing has occurred.
  
  \item \textbf{Advanced Glenohumeral Osteoarthritis or Severe Rotator Cuff Arthropathy:} These conditions may limit the potential for motion gain from capsular release alone, and alternative procedures (e.g., shoulder arthroplasty) might be more appropriate.
  
  \item \textbf{Neurological Deficits:} Pre-existing neurological deficits in the affected limb may complicate postoperative assessment of nerve function and rehabilitation.
  
  \item \textbf{Diabetes Mellitus:} While not a contraindication, diabetic patients often have a more severe and refractory course of adhesive capsulitis and may have a higher risk of recurrence or incomplete motion gain. Meticulous glycemic control is essential.
  
  \item \textbf{Excessive Force During Manipulation:} Extreme caution is paramount during manipulation under anesthesia to prevent iatrogenic injuries such as humeral fracture, rotator cuff tear, or brachial plexus/axillary nerve injury.
\end{itemize}

\section*{Complications \& Risks}
As with any surgical procedure, frozen shoulder release carries potential risks and complications, which can be broadly categorized into general surgical risks and procedure-specific risks.

\textbf{General Surgical Complications:}
\begin{itemize}
  \item \textbf{Anesthesia-related risks:} Adverse reactions to anesthetic agents, respiratory complications, cardiovascular events.
  
  \item \textbf{Infection:} Surgical site infection (superficial or deep), though rare with prophylactic antibiotics and sterile technique.
  
  \item \textbf{Bleeding and Hematoma:} Intraoperative or postoperative bleeding, potentially leading to hematoma formation.
  
  \item \textbf{Thromboembolic events:} Deep vein thrombosis (DVT) or pulmonary embolism (PE), though uncommon in shoulder surgery.
  
  \item \textbf{Pain:} Persistent or worsening pain despite the release.
\end{itemize}

\textbf{Procedure-Specific Complications \& Risks:}
\begin{itemize}
  \item \textbf{Nerve Injury:} The most commonly injured nerves are the axillary nerve (due to its proximity to the inferior capsule) and, less commonly, the musculocutaneous or suprascapular nerves. This can result in weakness or numbness.
  
  \item \textbf{Humeral Fracture:} A rare but serious complication, primarily associated with forceful manipulation under anesthesia, especially in patients with osteoporosis.
  
  \item \textbf{Recurrence of Stiffness:} If postoperative physiotherapy is inadequate or if the underlying inflammatory process is particularly aggressive, stiffness can recur, necessitating further intervention.
  
  \item \textbf{Incomplete Restoration of Motion:} Despite a thorough release, some patients may not regain full range of motion, particularly in severe, long-standing cases or in patients with comorbidities like diabetes.
  
  \item \textbf{Chondral Injury:} Damage to the articular cartilage of the humeral head or glenoid during instrument insertion or capsular release.
  
  \item \textbf{Rotator Cuff Injury:} Iatrogenic tear of a rotator cuff tendon, particularly during aggressive manipulation or if pre-existing degeneration is present.
  
  \item \textbf{Reflex Sympathetic Dystrophy (RSD) / Complex Regional Pain Syndrome (CRPS):} A rare but debilitating chronic pain condition that can develop after trauma or surgery.
  
  \item \textbf{Vascular Injury:} Damage to major blood vessels (e.g., axillary artery) is extremely rare but potentially severe.
\end{itemize}

\section*{Postoperative Recovery Timeline}
The postoperative recovery timeline for frozen shoulder release is highly dependent on patient compliance with an intensive rehabilitation program.

Immediately after surgery, patients are transferred to the Post-Anesthesia Care Unit (PACU) for monitoring of vital signs, pain levels, and neurovascular status of the affected limb. The regional nerve block typically provides several hours of profound pain relief. Patients are encouraged to begin gentle passive and active-assisted range-of-motion exercises as soon as tolerated, often under the guidance of a physical therapist. A sling may be used for comfort only, but prolonged immobilization is discouraged to prevent recurrence of stiffness. Pain is managed with a combination of oral analgesics, anti-inflammatory medications, and ice application to reduce swelling.

Over the first 1 to 2 weeks, the focus is on aggressive, supervised physiotherapy, often several times per week, complemented by a diligent home exercise program. The goal is to maintain and progressively increase the range of motion gained during surgery. Patients are typically restricted from heavy lifting or strenuous activities but can perform light activities of daily living. Wound care involves keeping the portal sites clean and dry until sutures or steri-strips are removed, usually around 7-14 days post-op. Long-term recovery extends over several weeks to months, with gradual progression of strengthening exercises and return to full activities. Full functional recovery, including near-normal range of motion and strength, can take 3 to 6 months, with continued improvement possible for up to a year. Consistent adherence to the rehabilitation protocol is the single most important factor influencing the speed and completeness of recovery.

\section*{Surgical Findings \& Pathology}
During frozen shoulder release, the surgeon typically observes characteristic findings within the glenohumeral joint. Arthroscopic examination reveals a globally contracted joint capsule, which appears thickened, fibrotic, and often inflamed. The joint volume is markedly reduced, making instrument maneuverability challenging. A hallmark finding is the obliteration of the rotator interval, the triangular space between the supraspinatus and subscapularis tendons, which is normally patent. The coracohumeral ligament is often visibly thickened and taut. The anterior and inferior capsular folds, which normally allow for significant redundancy during shoulder movement, are found to be severely contracted and adherent, restricting external rotation and abduction. Synovial tissue may appear hyperemic and hypertrophic, indicative of chronic inflammation.

Pathological examination of resected capsular tissue, if sent for analysis, typically confirms these macroscopic findings. Histologically, the tissue demonstrates features of chronic inflammation, characterized by an infiltration of inflammatory cells, synovial hyperplasia, and significant fibrosis. There is an increase in collagen deposition and myofibroblast proliferation within the capsular matrix, leading to the observed thickening and contracture. These findings underscore the fibrotic and inflammatory nature of adhesive capsulitis, which is mechanically addressed by surgical release.

\section*{Expected Postoperative Course}
A normal and successful postoperative course following frozen shoulder release is characterized by an immediate and significant improvement in passive range of motion, particularly in external rotation and abduction, which were often the most restricted movements preoperatively. While some pain is expected in the initial days, it should be manageable with oral analgesics and gradually subside. The key to a successful outcome is the diligent and consistent adherence to the prescribed physiotherapy regimen and home exercise program, which begins almost immediately after surgery.

Patients should expect a progressive increase in both passive and active range of motion over the ensuing weeks and months. The initial gains from surgery are mechanical, but sustained rehabilitation is crucial to prevent the capsule from scarring down again. Pain levels should steadily decrease as motion improves and inflammation resolves. Most patients will experience a gradual return to functional activities, including reaching, dressing, and sleeping comfortably, within 6 to 12 weeks. Full recovery, encompassing near-normal strength and endurance, may take several months, but the ultimate goal is a functional, mobile, and pain-free shoulder that allows the patient to resume their desired activities without significant limitation.

\section*{Postoperative Care \& Follow-up}
Postoperative care for frozen shoulder release is centered on intensive rehabilitation and close monitoring to ensure sustained motion gains and prevent recurrence.

\begin{itemize}
  \item \textbf{Physiotherapy:} This is the cornerstone of postoperative care. Patients begin a structured, aggressive physiotherapy program immediately after surgery, often within 24-48 hours. This typically involves daily home exercises and supervised sessions several times per week for the first 6-12 weeks, gradually progressing from passive and active-assisted range of motion to active motion and strengthening exercises.
  
  \item \textbf{Medication Management:} Pain is managed with oral analgesics (e.g., acetaminophen, NSAIDs, short-course opioids if necessary). Anti-inflammatory medications are often continued to help manage residual inflammation.
  
  \item \textbf{Wound Care:} The small arthroscopic portal incisions are typically covered with sterile dressings. Sutures or adhesive strips are usually removed by the surgeon or a nurse at the first follow-up visit, typically 7-14 days post-op. Patients are advised to keep the wounds clean and dry.
  
  \item \textbf{Follow-up Visits:} Regular follow-up appointments with the orthopedic surgeon are scheduled at intervals (e.g., 2 weeks, 6 weeks, 3 months, 6 months) to assess range of motion, pain levels, functional progress, and to adjust the rehabilitation program as needed.
  
  \item \textbf{Activity Restrictions:} While early motion is encouraged, patients are typically advised to avoid heavy lifting, pushing, or pulling for several weeks to allow for soft tissue healing and to prevent undue stress on the healing capsule. Gradual return to full activities, including sports, is guided by the surgeon and physical therapist.
  
  \item \textbf{Warning Signs:} Patients are educated on warning signs that warrant immediate contact with the surgeon, including fever, chills, worsening or uncontrolled pain, increased redness or swelling around the surgical sites, purulent discharge from the wounds, new or worsening numbness or weakness in the arm or hand, or a sudden loss of the motion gained after surgery.
\end{itemize}

Long-term compliance with a home exercise program is crucial even after formal physiotherapy ends to maintain flexibility and prevent recurrence.

\section*{Epidemiology}
Adhesive capsulitis, or frozen shoulder, is a relatively common condition affecting approximately 2\% to 5\% of the general population. Its incidence peaks in individuals between 40 and 60 years of age, with a slight predilection for women. While it can affect either shoulder, about 20\% to 30\% of patients will eventually develop the condition in the contralateral shoulder, often within 5 years of the initial presentation.

A significant epidemiological association exists with certain systemic diseases. Diabetes mellitus is the most prominent risk factor, increasing the incidence of frozen shoulder by two to four times; up to 20\% of diabetic patients may experience adhesive capsulitis. Diabetics also tend to have a more severe, prolonged, and refractory course, with a higher likelihood of requiring surgical intervention and a greater risk of incomplete recovery. Other associated conditions include thyroid disease (hypothyroidism and hyperthyroidism), cardiovascular disease, Parkinson's disease, and certain medications (e.g., antiretroviral therapy, phenobarbital). Previous shoulder trauma, surgery, or prolonged immobilization can also precipitate secondary adhesive capsulitis. While the condition is generally self-limiting over 1-2 years, the morbidity associated with severe pain and functional limitation can be substantial, impacting quality of life and productivity. Mortality directly related to frozen shoulder is negligible, but the morbidity from chronic pain and disability can be significant.

\section*{Outcomes \& Prognosis}
The prognosis for frozen shoulder, even without surgical intervention, is generally favorable, with most patients experiencing gradual improvement over 1 to 2 years. However, up to one-third of patients may be left with some degree of permanent stiffness or pain. For appropriately selected patients who have failed conservative management, surgical release offers a more rapid and predictable restoration of motion and pain relief.

Numerous studies and systematic reviews consistently report high success rates for arthroscopic capsular release, with significant improvements in both passive and active range of motion across all planes. Patient satisfaction rates typically range from 80\% to 95\%. Functional outcomes, including the ability to perform activities of daily living, return to work, and participate in recreational sports, are generally excellent. The risk of recurrence after a comprehensive surgical release followed by an intensive rehabilitation program is relatively low, estimated to be less than 10\%. Prognostic factors influencing outcomes include the duration and severity of symptoms prior to surgery, the presence of comorbidities such as diabetes (which can lead to a more prolonged and less complete recovery), and, most critically, patient adherence to the postoperative physiotherapy regimen. While there isn't a single "landmark trial" for frozen shoulder release akin to those for other surgical conditions, the cumulative evidence from numerous cohort studies and meta-analyses supports its efficacy as a safe and effective treatment for refractory adhesive capsulitis.

\section*{References}
\begin{enumerate}
  \item MedlinePlus. Frozen shoulder surgery. U.S. National Library of Medicine; 2024.
  
  \item Canale ST, Beaty JH. Campbell's Operative Orthopaedics, 14th ed. Elsevier; 2021.
  
  \item Rockwood CA Jr, Matsen FA III, Wirth MA, et al. Rockwood and Matsen's The Shoulder, 6th ed. Elsevier; 2021.
  
  \item American Academy of Orthopaedic Surgeons. Management of adhesive capsulitis of the shoulder: evidence-based clinical practice guideline. AAOS; 2021.
\end{enumerate}

\clearpage